% ===========================================================================
%  evnx for AI / ML Engineers  —  Python, Jupyter, GPU boxes, Rust inference
%  Persona SM.  Build: pdflatex -interaction=nonstopmode ai-ml.tex
% ===========================================================================
\documentclass[aspectratio=169,10pt]{beamer}

\newcommand{\capturedir}{../captures}
% ---------------------------------------------------------------------------
%  Shared preamble for the seven audience decks.
%
%  Each deck does, before the document body starts:
%     \newcommand{\capturedir}{../captures}
%     % ---------------------------------------------------------------------------
%  Shared preamble for the seven audience decks.
%
%  Each deck does, before the document body starts:
%     \newcommand{\capturedir}{../captures}
%     % ---------------------------------------------------------------------------
%  Shared preamble for the seven audience decks.
%
%  Each deck does, before the document body starts:
%     \newcommand{\capturedir}{../captures}
%     \input{persona-preamble}
%
%  \capturedir points at ../captures so both pools resolve:
%     \capture{24-scan.txt}              — the original v0.5.0 captures
%     \capture{persona/fe-01-scan-dist.txt} — captures recorded for these decks
% ---------------------------------------------------------------------------

\input{../common/evnx-preamble}

% ---- provenance ------------------------------------------------------------
% The original decks were pinned to tag v0.5.0 (85a52c3). Everything captured
% for the persona decks was re-run against 0.5.2, so the footline differs.
\renewcommand{\verifiedon}{\tiny\textcolor{evnxMuted}{verified against evnx 0.5.2 --- real run, not transcribed}}

% ---- audience badge on the title slide ------------------------------------
\newcommand{\audience}[1]{%
  \begin{center}\small\textcolor{evnxAccent}{\textbf{Audience: #1}}\end{center}}

% ---- verdict chips ---------------------------------------------------------
\newcommand{\solved}{\textcolor{evnxOk}{\ensuremath{\checkmark}}}
\newcommand{\partialmark}{\textcolor{evnxWarn}{\textbf{\textasciitilde}}}
\newcommand{\unsolved}{\textcolor{evnxErr}{\ensuremath{\times}}}

% ---- a correction block: something the source material got wrong -----------
\newenvironment{correction}[1]{%
  \begin{alertblock}{\small Correction --- #1}\small}{\end{alertblock}}

% ---- a trap block: works in the demo, fails in real use --------------------
\newenvironment{trap}[1]{%
  \setbeamercolor{block title}{bg=evnxWarn,fg=white}%
  \setbeamercolor{block body}{bg=evnxWarn!8}%
  \begin{block}{\small #1}\small}{\end{block}}

% ---- the persona card ------------------------------------------------------
\newcommand{\personacard}[4]{%
  \begin{block}{#1}
    \small
    \textbf{Stack:} #2 \\
    \textbf{Ships:} #3 \\
    \textbf{Loses time to:} #4
  \end{block}}

% ---- speaker note (rendered small and muted, not hidden) -------------------
\newcommand{\saythis}[1]{%
  \vspace{0.4em}{\footnotesize\textcolor{evnxMuted}{\textit{Say:} #1}}}

%
%  \capturedir points at ../captures so both pools resolve:
%     \capture{24-scan.txt}              — the original v0.5.0 captures
%     \capture{persona/fe-01-scan-dist.txt} — captures recorded for these decks
% ---------------------------------------------------------------------------

% ---------------------------------------------------------------------------
%  evnx presentation preamble — shared by all three decks
%
%  Engine: pdfLaTeX (also builds under XeLaTeX / LuaLaTeX).
%  Packages: beamer, listings, xcolor, amssymb, newunicodechar, booktabs,
%            fancyvrb, hyperref — all in texlive-latex-recommended + beamer.
%
%  \capturedir must be defined by the including document BEFORE \input of
%  this file, e.g.  \newcommand{\capturedir}{../captures}
% ---------------------------------------------------------------------------

\usetheme{Madrid}
\usecolortheme{seahorse}
\usefonttheme{professionalfonts}
\setbeamertemplate{navigation symbols}{}
\setbeamertemplate{caption}[numbered]
\setbeamerfont{frametitle}{size=\large}
\setbeamerfont{title}{size=\LARGE,series=\bfseries}

\usepackage[T1]{fontenc}
\usepackage[utf8]{inputenc}
\usepackage{lmodern}
\usepackage{amssymb}
\usepackage{booktabs}
\usepackage{xcolor}
\usepackage{listings}
\usepackage{fancyvrb}
\usepackage{newunicodechar}
\usepackage{tikz}
\usetikzlibrary{arrows.meta,positioning,shapes.geometric,fit,backgrounds}

% ---- palette ---------------------------------------------------------------
\definecolor{evnxInk}{HTML}{1F2933}
\definecolor{evnxAccent}{HTML}{2F6F4E}
\definecolor{evnxWarn}{HTML}{B7791F}
\definecolor{evnxErr}{HTML}{B03A2E}
\definecolor{evnxOk}{HTML}{2F6F4E}
\definecolor{evnxMuted}{HTML}{6B7280}
\definecolor{evnxTermBg}{HTML}{F5F5F2}
\definecolor{evnxRule}{HTML}{D1D5DB}

\setbeamercolor{structure}{fg=evnxAccent}
\setbeamercolor{alerted text}{fg=evnxErr}

% ---- glyph mapping for body text ------------------------------------------
% evnx emits exactly these 11 non-ASCII characters (enumerated from the
% captured logs, not guessed).
\newcommand{\evnxCheck}{\textcolor{evnxOk}{\ensuremath{\checkmark}}}
\newcommand{\evnxCross}{\textcolor{evnxErr}{\ensuremath{\times}}}
\newcommand{\evnxWarnSign}{\textcolor{evnxWarn}{\textbf{!}}}

\newunicodechar{✓}{\evnxCheck}
\newunicodechar{✗}{\evnxCross}
\newunicodechar{⚠}{\evnxWarnSign}
\newunicodechar{→}{\ensuremath{\rightarrow}}
\newunicodechar{↔}{\ensuremath{\leftrightarrow}}
\newunicodechar{─}{\textendash}
\newunicodechar{—}{---}
\newunicodechar{·}{\textperiodcentered}
\newunicodechar{…}{\ldots}
\newunicodechar{•}{\textbullet}
\newunicodechar{️}{}                % U+FE0F variation selector — invisible

% ---- terminal listing style ------------------------------------------------
% `literate` is required: newunicodechar does not apply inside listings.
\lstdefinestyle{evnxterm}{
  basicstyle=\ttfamily\scriptsize,
  backgroundcolor=\color{evnxTermBg},
  frame=single,
  rulecolor=\color{evnxRule},
  framesep=4pt,
  breaklines=true,
  breakatwhitespace=false,
  columns=fullflexible,
  keepspaces=true,
  showstringspaces=false,
  upquote=true,
  extendedchars=true,
  literate=
    {✓}{{\textcolor{evnxOk}{$\checkmark$}}}1
    {✗}{{\textcolor{evnxErr}{$\times$}}}1
    {⚠}{{\textcolor{evnxWarn}{\textbf{!}}}}1
    {→}{{$\rightarrow$}}1
    {↔}{{$\leftrightarrow$}}1
    {─}{{-}}1
    {—}{{---}}1
    {·}{{$\cdot$}}1
    {…}{{\ldots}}1
    {•}{{\textbullet}}1
    {️}{{}}0
    {é}{{\'e}}1
}

\lstdefinestyle{evnxcode}{
  style=evnxterm,
  basicstyle=\ttfamily\scriptsize,
  backgroundcolor=\color{white},
}

\lstset{style=evnxterm}

% ---- helper macros ---------------------------------------------------------
% \capture{file}        — include a real captured run verbatim
% \capturepart{f}{a}{b} — include lines a..b of a captured run
\newcommand{\capture}[1]{\lstinputlisting[style=evnxterm]{\capturedir/#1}}
\newcommand{\capturepart}[3]{%
  \lstinputlisting[style=evnxterm,firstline=#2,lastline=#3]{\capturedir/#1}}

% a labelled "real output" box
\newcommand{\realrun}[1]{%
  \begin{center}\scriptsize\textcolor{evnxMuted}{real run \textendash\ #1}\end{center}}

\newcommand{\cmd}[1]{\texttt{\textbf{#1}}}
\newcommand{\flag}[1]{\texttt{#1}}
\newcommand{\file}[1]{\texttt{#1}}

% exit-code chips
\newcommand{\exitzero}{\textcolor{evnxOk}{\texttt{0}}}
\newcommand{\exitone}{\textcolor{evnxWarn}{\texttt{1}}}
\newcommand{\exittwo}{\textcolor{evnxErr}{\texttt{2}}}

% a standard "verified against" footline for factual slides
\newcommand{\verifiedon}{\tiny\textcolor{evnxMuted}{verified against evnx 0.5.0 @ 9cfe75b}}

\usepackage[colorlinks=true,linkcolor=evnxAccent,urlcolor=evnxAccent]{hyperref}


% ---- provenance ------------------------------------------------------------
% The original decks were pinned to tag v0.5.0 (85a52c3). Everything captured
% for the persona decks was re-run against 0.5.2, so the footline differs.
\renewcommand{\verifiedon}{\tiny\textcolor{evnxMuted}{verified against evnx 0.5.2 --- real run, not transcribed}}

% ---- audience badge on the title slide ------------------------------------
\newcommand{\audience}[1]{%
  \begin{center}\small\textcolor{evnxAccent}{\textbf{Audience: #1}}\end{center}}

% ---- verdict chips ---------------------------------------------------------
\newcommand{\solved}{\textcolor{evnxOk}{\ensuremath{\checkmark}}}
\newcommand{\partialmark}{\textcolor{evnxWarn}{\textbf{\textasciitilde}}}
\newcommand{\unsolved}{\textcolor{evnxErr}{\ensuremath{\times}}}

% ---- a correction block: something the source material got wrong -----------
\newenvironment{correction}[1]{%
  \begin{alertblock}{\small Correction --- #1}\small}{\end{alertblock}}

% ---- a trap block: works in the demo, fails in real use --------------------
\newenvironment{trap}[1]{%
  \setbeamercolor{block title}{bg=evnxWarn,fg=white}%
  \setbeamercolor{block body}{bg=evnxWarn!8}%
  \begin{block}{\small #1}\small}{\end{block}}

% ---- the persona card ------------------------------------------------------
\newcommand{\personacard}[4]{%
  \begin{block}{#1}
    \small
    \textbf{Stack:} #2 \\
    \textbf{Ships:} #3 \\
    \textbf{Loses time to:} #4
  \end{block}}

% ---- speaker note (rendered small and muted, not hidden) -------------------
\newcommand{\saythis}[1]{%
  \vspace{0.4em}{\footnotesize\textcolor{evnxMuted}{\textit{Say:} #1}}}

%
%  \capturedir points at ../captures so both pools resolve:
%     \capture{24-scan.txt}              — the original v0.5.0 captures
%     \capture{persona/fe-01-scan-dist.txt} — captures recorded for these decks
% ---------------------------------------------------------------------------

% ---------------------------------------------------------------------------
%  evnx presentation preamble — shared by all three decks
%
%  Engine: pdfLaTeX (also builds under XeLaTeX / LuaLaTeX).
%  Packages: beamer, listings, xcolor, amssymb, newunicodechar, booktabs,
%            fancyvrb, hyperref — all in texlive-latex-recommended + beamer.
%
%  \capturedir must be defined by the including document BEFORE \input of
%  this file, e.g.  \newcommand{\capturedir}{../captures}
% ---------------------------------------------------------------------------

\usetheme{Madrid}
\usecolortheme{seahorse}
\usefonttheme{professionalfonts}
\setbeamertemplate{navigation symbols}{}
\setbeamertemplate{caption}[numbered]
\setbeamerfont{frametitle}{size=\large}
\setbeamerfont{title}{size=\LARGE,series=\bfseries}

\usepackage[T1]{fontenc}
\usepackage[utf8]{inputenc}
\usepackage{lmodern}
\usepackage{amssymb}
\usepackage{booktabs}
\usepackage{xcolor}
\usepackage{listings}
\usepackage{fancyvrb}
\usepackage{newunicodechar}
\usepackage{tikz}
\usetikzlibrary{arrows.meta,positioning,shapes.geometric,fit,backgrounds}

% ---- palette ---------------------------------------------------------------
\definecolor{evnxInk}{HTML}{1F2933}
\definecolor{evnxAccent}{HTML}{2F6F4E}
\definecolor{evnxWarn}{HTML}{B7791F}
\definecolor{evnxErr}{HTML}{B03A2E}
\definecolor{evnxOk}{HTML}{2F6F4E}
\definecolor{evnxMuted}{HTML}{6B7280}
\definecolor{evnxTermBg}{HTML}{F5F5F2}
\definecolor{evnxRule}{HTML}{D1D5DB}

\setbeamercolor{structure}{fg=evnxAccent}
\setbeamercolor{alerted text}{fg=evnxErr}

% ---- glyph mapping for body text ------------------------------------------
% evnx emits exactly these 11 non-ASCII characters (enumerated from the
% captured logs, not guessed).
\newcommand{\evnxCheck}{\textcolor{evnxOk}{\ensuremath{\checkmark}}}
\newcommand{\evnxCross}{\textcolor{evnxErr}{\ensuremath{\times}}}
\newcommand{\evnxWarnSign}{\textcolor{evnxWarn}{\textbf{!}}}

\newunicodechar{✓}{\evnxCheck}
\newunicodechar{✗}{\evnxCross}
\newunicodechar{⚠}{\evnxWarnSign}
\newunicodechar{→}{\ensuremath{\rightarrow}}
\newunicodechar{↔}{\ensuremath{\leftrightarrow}}
\newunicodechar{─}{\textendash}
\newunicodechar{—}{---}
\newunicodechar{·}{\textperiodcentered}
\newunicodechar{…}{\ldots}
\newunicodechar{•}{\textbullet}
\newunicodechar{️}{}                % U+FE0F variation selector — invisible

% ---- terminal listing style ------------------------------------------------
% `literate` is required: newunicodechar does not apply inside listings.
\lstdefinestyle{evnxterm}{
  basicstyle=\ttfamily\scriptsize,
  backgroundcolor=\color{evnxTermBg},
  frame=single,
  rulecolor=\color{evnxRule},
  framesep=4pt,
  breaklines=true,
  breakatwhitespace=false,
  columns=fullflexible,
  keepspaces=true,
  showstringspaces=false,
  upquote=true,
  extendedchars=true,
  literate=
    {✓}{{\textcolor{evnxOk}{$\checkmark$}}}1
    {✗}{{\textcolor{evnxErr}{$\times$}}}1
    {⚠}{{\textcolor{evnxWarn}{\textbf{!}}}}1
    {→}{{$\rightarrow$}}1
    {↔}{{$\leftrightarrow$}}1
    {─}{{-}}1
    {—}{{---}}1
    {·}{{$\cdot$}}1
    {…}{{\ldots}}1
    {•}{{\textbullet}}1
    {️}{{}}0
    {é}{{\'e}}1
}

\lstdefinestyle{evnxcode}{
  style=evnxterm,
  basicstyle=\ttfamily\scriptsize,
  backgroundcolor=\color{white},
}

\lstset{style=evnxterm}

% ---- helper macros ---------------------------------------------------------
% \capture{file}        — include a real captured run verbatim
% \capturepart{f}{a}{b} — include lines a..b of a captured run
\newcommand{\capture}[1]{\lstinputlisting[style=evnxterm]{\capturedir/#1}}
\newcommand{\capturepart}[3]{%
  \lstinputlisting[style=evnxterm,firstline=#2,lastline=#3]{\capturedir/#1}}

% a labelled "real output" box
\newcommand{\realrun}[1]{%
  \begin{center}\scriptsize\textcolor{evnxMuted}{real run \textendash\ #1}\end{center}}

\newcommand{\cmd}[1]{\texttt{\textbf{#1}}}
\newcommand{\flag}[1]{\texttt{#1}}
\newcommand{\file}[1]{\texttt{#1}}

% exit-code chips
\newcommand{\exitzero}{\textcolor{evnxOk}{\texttt{0}}}
\newcommand{\exitone}{\textcolor{evnxWarn}{\texttt{1}}}
\newcommand{\exittwo}{\textcolor{evnxErr}{\texttt{2}}}

% a standard "verified against" footline for factual slides
\newcommand{\verifiedon}{\tiny\textcolor{evnxMuted}{verified against evnx 0.5.0 @ 9cfe75b}}

\usepackage[colorlinks=true,linkcolor=evnxAccent,urlcolor=evnxAccent]{hyperref}


% ---- provenance ------------------------------------------------------------
% The original decks were pinned to tag v0.5.0 (85a52c3). Everything captured
% for the persona decks was re-run against 0.5.2, so the footline differs.
\renewcommand{\verifiedon}{\tiny\textcolor{evnxMuted}{verified against evnx 0.5.2 --- real run, not transcribed}}

% ---- audience badge on the title slide ------------------------------------
\newcommand{\audience}[1]{%
  \begin{center}\small\textcolor{evnxAccent}{\textbf{Audience: #1}}\end{center}}

% ---- verdict chips ---------------------------------------------------------
\newcommand{\solved}{\textcolor{evnxOk}{\ensuremath{\checkmark}}}
\newcommand{\partialmark}{\textcolor{evnxWarn}{\textbf{\textasciitilde}}}
\newcommand{\unsolved}{\textcolor{evnxErr}{\ensuremath{\times}}}

% ---- a correction block: something the source material got wrong -----------
\newenvironment{correction}[1]{%
  \begin{alertblock}{\small Correction --- #1}\small}{\end{alertblock}}

% ---- a trap block: works in the demo, fails in real use --------------------
\newenvironment{trap}[1]{%
  \setbeamercolor{block title}{bg=evnxWarn,fg=white}%
  \setbeamercolor{block body}{bg=evnxWarn!8}%
  \begin{block}{\small #1}\small}{\end{block}}

% ---- the persona card ------------------------------------------------------
\newcommand{\personacard}[4]{%
  \begin{block}{#1}
    \small
    \textbf{Stack:} #2 \\
    \textbf{Ships:} #3 \\
    \textbf{Loses time to:} #4
  \end{block}}

% ---- speaker note (rendered small and muted, not hidden) -------------------
\newcommand{\saythis}[1]{%
  \vspace{0.4em}{\footnotesize\textcolor{evnxMuted}{\textit{Say:} #1}}}


\title{evnx for AI and ML engineers}
\subtitle{Expensive keys, disposable machines, and notebooks that remember}
\date{}

\begin{document}

\begin{frame}[plain]
  \titlepage
  \audience{Python · Jupyter · LLM APIs · rented GPUs · Slurm · Rust inference}
  \begin{center}\footnotesize
  \textcolor{evnxMuted}{evnx 0.5.2. Every terminal block is a real run,
  including the two where evnx misses.}
  \end{center}
\end{frame}

% =========================================================================
\begin{frame}{This deck in one slide}
  \begin{block}{Your threat model is different from everyone else's}
    A leaked web API key is an incident. A leaked \textbf{LLM} key is an
    incident \emph{and} an invoice. And your secrets live on machines you
    rent by the hour and destroy by lunchtime.
  \end{block}

  \vspace{0.6em}
  \begin{block}{Two things you must know before you trust \cmd{scan}}
    \begin{enumerate}
      \item \textbf{\file{.ipynb} files are not scanned at all.} Not
            low-confidence --- not scanned.
      \item The OpenAI detector matches the \textbf{legacy} key format.
            Modern \texttt{sk-proj-} keys fall through to generic entropy at
            \textbf{low} severity, which \flag{-{}-severity high} filters out.
    \end{enumerate}
  \end{block}
  \saythis{Both are demoed. Do not soften them --- this audience will test it.}
\end{frame}

\begin{frame}{Contents}\tableofcontents[hideallsubsections]\end{frame}

% =========================================================================
\section{Three things that already happened to you}

\begin{frame}{A month in the life}
  \small
  \begin{block}{The notebook}
    You paste \texttt{OPENAI\_API\_KEY="sk-..."} into a cell ``just to test'',
    run it, commit the notebook --- \textbf{with the output cell} still showing
    the key. Three days later the provider emails about unusual usage.
  \end{block}

  \begin{block}{The GPU boxes}
    You rent two machines for a fine-tune and \texttt{scp} your laptop
    \file{.env} to both, including production AWS keys you did not need. The
    machines are destroyed --- but a copy lives on in a teammate's home
    directory on the Slurm cluster.
  \end{block}

  \begin{block}{The intern}
    Six weeks, needs the vector DB and W\&B keys. You share your personal keys
    over chat. When they leave, nobody rotates anything.
  \end{block}
\end{frame}

\begin{frame}{What evnx changes, per incident}
  \small
  \begin{tabular}{@{}lll@{}}
    \toprule
    \textbf{Incident} & \textbf{Mechanism} & \textbf{Verdict} \\
    \midrule
    Notebook commit  & \cmd{scan} in pre-commit & \unsolved\ for \file{.ipynb} --- see slide 8 \\
    GPU box copies   & \cmd{cloud run} --- no file written & \solved \\
    Intern access    & dedicated vault + \cmd{vault revoke} & \partialmark\ no expiry \\
    \bottomrule
  \end{tabular}

  \vspace{1em}
  \begin{block}{The one that generalises}
    \cmd{evnx cloud run -{}- python train.py} means \textbf{there is no
    \file{.env} to leave behind} on a machine you are about to destroy, or on a
    shared cluster filesystem. That is the single highest-value change for this
    audience and it works today.
  \end{block}
\end{frame}

% =========================================================================
\section{The notebook problem}

\begin{frame}[fragile]{Setup: one key, two identical files}
\begin{lstlisting}[style=evnxcode]
notebooks/explore.ipynb    # a real notebook with the key in a code cell
notebooks/explore.json     # byte-identical copy, different extension
\end{lstlisting}

  \vspace{0.6em}
  Same bytes. Same key. The only difference is the file extension.

  \vspace{0.8em}
  \saythis{``What should a secret scanner do here?'' --- let them say ``flag both''.}
\end{frame}

\begin{frame}[fragile]{\cmd{evnx scan notebooks/}}
  \capture{persona/ml-01-scan-notebooks.txt}
  \realrun{two identical files, one finding}

  \vspace{0.4em}
  \begin{trap}{Read the ``scanning'' line}
    It scanned \file{explore.json}. It never opened \file{explore.ipynb}.
    \texttt{ipynb} is missing from the scannable-extension allowlist at
    \file{src/commands/scan/filters.rs:284}. Tracked as \textbf{NEW-2}.
  \end{trap}
\end{frame}

\begin{frame}[fragile]{And naming the notebook directly does not help}
  \capture{persona/ml-02-scan-ipynb.txt}
  \realrun{the notebook, named explicitly --- exit \exitzero}

  \vspace{0.4em}
  {\footnotesize ``0 files scanned'' and a green check. If your pre-commit hook
  is \cmd{evnx scan notebooks/}, it has been passing without doing anything.}

  \vspace{0.5em}
  \begin{block}{Workaround, today}
    \begin{center}\ttfamily\scriptsize
    find . -name '*.ipynb' -exec sh -c 'cp "\$1" "\$1.json"' \_ \{\} \textbackslash; \&\& evnx scan . ; \\
    find . -name '*.ipynb.json' -delete
    \end{center}
    \vspace{0.3em}
    Ugly, but it is the difference between a scanner and a green tick.
  \end{block}
\end{frame}

\begin{frame}{Do this regardless of which scanner you use}
  \begin{center}\large\ttfamily
  pip install nbstripout \&\& nbstripout -{}-install
  \end{center}

  \vspace{1em}
  \begin{block}{Why this is first, not fifth}
    The notebook incident is not really ``a key in source''. It is
    \textbf{a key in a saved output cell}. Stripping outputs removes the whole
    class --- and it works whether or not any scanner understands
    \file{.ipynb}.
  \end{block}

  \vspace{0.8em}
  \saythis{Recommending the non-evnx tool here is what buys you credibility for
  the rest of the deck.}
\end{frame}

% =========================================================================
\section{The detector problem}

\begin{frame}[fragile]{Two OpenAI keys, one file}
  \capturepart{persona/ml-06-openai-key-formats.txt}{1}{18}
  \realrun{line 1 is a legacy \texttt{sk-} key, line 2 is a \texttt{sk-proj-} key}
\end{frame}

\begin{frame}{Why one is \texttt{high} and the other is \texttt{low}}
  \begin{center}\ttfamily\small
  pattern: r"sk-[0-9a-zA-Z]\{48\}"
  \end{center}
  \begin{center}\footnotesize\textcolor{evnxMuted}{src/utils/patterns.rs:106}\end{center}

  \vspace{0.8em}
  A modern project key is \texttt{sk-proj-} followed by a mix that includes
  \texttt{-} and \texttt{\_}. The character class is alphanumeric only, so the
  hyphen after \texttt{proj} ends the match. The key is then caught only by the
  generic Shannon-entropy detector, at \textbf{low} confidence.

  \vspace{0.8em}
  \begin{trap}{The consequence for your CI}
    \cmd{evnx scan . -{}-severity high} --- the sensible setting for a blocking
    gate --- \textbf{will not report a modern OpenAI key}. We verified it:
    ``✓ No secrets detected'', exit \exitzero.
  \end{trap}
\end{frame}

\begin{frame}[fragile]{Proof, so nobody has to take my word for it}
  \capture{persona/ml-05-scan-severity-high.txt}
  \realrun{the same notebook directory, \flag{-{}-severity high}}

  \vspace{0.4em}
  \begin{block}{What to do this week}
    Add your own patterns rather than waiting:
    \vspace{0.3em}
    \begin{center}\ttfamily\scriptsize
    evnx scan . -{}-pattern 'sk-proj-[A-Za-z0-9\_-]\{40,\}' \textbackslash \\
    \ \ \ \ \ \ \ \ \ \ \ \ \ -{}-pattern 'hf\_[A-Za-z0-9]\{30,\}' \textbackslash \\
    \ \ \ \ \ \ \ \ \ \ \ \ \ -{}-pattern 'r8\_[A-Za-z0-9]\{30,\}'
    \end{center}
  \end{block}
\end{frame}

\begin{frame}[fragile]{The other half: name-based detection does work}
  \capturepart{persona/ml-03-scan-env.txt}{1}{16}
  \realrun{\texttt{evnx scan .env}}

  \vspace{0.3em}
  {\footnotesize Note the annotation: \emph{``flagged by variable name --- the
  value was not checked''}. In a \file{.env} the key \emph{name} carries the
  signal, so \texttt{OPENAI\_API\_KEY} and \texttt{HF\_TOKEN} are both caught
  at \texttt{high} even though neither value matched a pattern. The gap is in
  \textbf{source files and notebooks}, where there is no variable name to lean on.}
\end{frame}

% =========================================================================
\section{Many machines, no files}

\begin{frame}[fragile]{The GPU box workflow}
\begin{lstlisting}[style=evnxcode]
# on the laptop, once
evnx vault create llm-research --env dev
evnx cloud link llm-research/dev && evnx cloud push

# on every rented box, every cluster node -- no scp, no file
evnx cloud run -- python train.py --config cfg/exp42.yaml
evnx cloud run -- jupyter lab
evnx cloud run -- streamlit run app.py
evnx cloud run -- cargo run          # the Rust inference service
\end{lstlisting}

  \vspace{0.6em}
  \begin{block}{Blast radius, by design}
    \begin{tabular}{@{}ll@{}}
      \texttt{llm-research/dev}  & personal keys, low spend limits \\
      \texttt{rag-api/prod}      & production keys, few members \\
      \texttt{demos/public}      & keys that are safe for interns and demos \\
    \end{tabular}
  \end{block}
  {\footnotesize Three vaults is the whole intern problem solved structurally ---
  they get \texttt{demos/public} and never see the other two.}
\end{frame}

\begin{frame}{Caveats worth stating before someone finds them}
  \small
  \begin{itemize}
    \item \textbf{\cmd{cloud run} still prompts for the master password} on the
          remote box. On a shared machine that is a real exposure --- machine
          identities are not built yet. Item 4 on the roadmap.
    \item \textbf{Jupyter kernels started elsewhere} do not inherit the
          environment. \cmd{cloud run -{}- jupyter lab} injects into the
          \emph{parent}; a kernel launched by a different supervisor will not
          see it.
    \item \textbf{Colab and hosted notebooks:} evnx cannot run there in any
          useful way. Use the platform's own secret store.
    \item \textbf{Slurm:} \cmd{cloud run -{}- srun python train.py} works from
          the login node, but whether the environment propagates depends on
          your cluster's configuration.
  \end{itemize}
\end{frame}

\begin{frame}[fragile]{The AI-assistant angle, which lands with this room}
  \begin{block}{The actual argument}
    A coding assistant with filesystem access reads your working tree. If
    there is no \file{.env} in it, there is nothing to read.
  \end{block}

  \vspace{0.6em}
  \begin{center}\large\cmd{evnx cloud run -{}- python app.py}\end{center}

  \vspace{0.8em}
  \begin{trap}{Do not oversell it}
    An assistant that can run \emph{commands} can run \cmd{evnx cloud pull}.
    This removes the passive read, not the active one.
  \end{trap}
\end{frame}

% =========================================================================
\section{Reproducibility}

\begin{frame}[fragile]{``Which keys and config did run \#42 use?''}
  \cmd{evnx cloud history} lists a vault's versions, newest first. It is
  \textbf{read-only}: no password, no decryption, no blob download --- so it is
  safe to call from inside a training script.

  \vspace{0.7em}
\begin{lstlisting}[style=evnxcode,language=Python]
import subprocess, wandb

hist = subprocess.check_output(
    ["evnx", "cloud", "history", "--limit", "1"], text=True)
wandb.config.update({"env_vault_head": hist.strip()})
\end{lstlisting}

  \vspace{0.5em}
  \begin{trap}{There is no \flag{-{}-format json} on \cmd{cloud status} or \cmd{cloud history}}
    Both are human-formatted only, so you are parsing text. A machine-readable
    status is item 6 on the roadmap --- until then, log the raw string rather
    than pretending to extract a field.
  \end{trap}
\end{frame}

\begin{frame}[fragile]{Custom detectors: the team-wide version}
  \capturepart{persona/ml-07-custom-pattern-toml.txt}{1}{20}
  \realrun{\texttt{[[scan.patterns]]} in \file{.evnx.toml}}

  \vspace{0.3em}
  {\footnotesize \flag{-{}-pattern} is for a one-off. \file{.evnx.toml} is
  committed, so the whole team scans with the same rules --- and both sets
  apply together. Note evnx names the unknown \texttt{severity} key and says it
  ignored it, rather than dropping it quietly.}
\end{frame}

\begin{frame}[fragile]{What validate is good for here}
  \capture{persona/ml-04-validate.txt}
  \realrun{a research \file{.env} with one unreplaced placeholder}

  \vspace{0.4em}
  {\footnotesize Catching \texttt{your\_wandb\_key\_here} at commit time is
  worth more than it sounds: the alternative is finding out forty minutes into
  a training run when the logger 401s.}
\end{frame}

% =========================================================================
\section{Adoption}

\begin{frame}[fragile]{Copy-paste}
\begin{lstlisting}[style=evnxcode]
# vaults by blast radius
evnx vault create llm-research --env dev
evnx vault create rag-api --env prod
evnx vault create demos --env public

# every machine: no plaintext file
evnx cloud run -- jupyter lab
evnx cloud run -- python train.py --config cfg/exp42.yaml

# guardrails
pip install nbstripout && nbstripout --install     # do this first
evnx scan . --pattern 'sk-proj-[A-Za-z0-9_-]{40,}' \
            --pattern 'hf_[A-Za-z0-9]{30,}'

# deploy the RAG API
evnx convert --to kubernetes    # or --to gcp-secrets / --to aws-secrets
\end{lstlisting}
\end{frame}

\begin{frame}{The scorecard, revised}
  \begin{center}
  \begin{tabular}{@{}clp{6.2cm}@{}}
    \toprule
    & \textbf{Count} & \textbf{Items} \\
    \midrule
    \solved     & 6 & many machines, vault-per-blast-radius, Streamlit/Gradio, leak prevention in \emph{source}, no-file-on-disk, K8s/GCP deploy \\
    \partialmark & 11 & experiment config, Jupyter kernels, placeholders, Rust binaries, credentialed URIs, rotation, intern access, Docker images, Slurm, run provenance, \textbf{file secrets} \\
    \unsolved   & 3 & hosted notebooks, HF Spaces, \textbf{\file{.ipynb} scanning} \\
    \bottomrule
  \end{tabular}
  \end{center}

  \vspace{0.5em}
  \begin{correction}{two moves from the source material}
    \file{.ipynb} scanning was listed \solved. It is \unsolved\ --- verified.
    \textbf{File secrets} (service-account JSON) were listed \unsolved; inline
    multi-line JSON \emph{does} parse if single-quoted, so it is \partialmark.
  \end{correction}
\end{frame}

\begin{frame}{What is coming, in priority order}
  \small
  \begin{tabular}{@{}clp{5.6cm}@{}}
    \toprule
    & \textbf{Feature} & \textbf{Fixes the slide on} \\
    \midrule
    P0 & Scan explicitly-named excluded/unscannable paths & ``0 files scanned'' \\
    1  & AI detector pack (HF, W\&B, Pinecone, LangSmith, Cohere, Replicate, Groq) + modern OpenAI format & \flag{-{}-severity high} missing keys \\
    2  & \file{.ipynb}-aware scanning (sources \emph{and} outputs) & the notebook \\
    3  & File secrets materialised to temp files by \cmd{cloud run} & \texttt{GOOGLE\_APPLICATION\_CREDENTIALS} \\
    4  & Machine tokens that decrypt one vault & GPU boxes, Slurm \\
    5  & HF Spaces / Modal / Replicate convert targets & demo deploys \\
    \bottomrule
  \end{tabular}
\end{frame}

\begin{frame}[plain]
  \vfill
  \begin{center}
    {\LARGE\textbf{Questions}}\\[1.2em]
    {\small \url{https://www.evnx.dev/guides} · \url{https://github.com/urwithajit9/evnx}}\\[0.8em]
    {\footnotesize\textcolor{evnxMuted}{evnx 0.5.2 · every terminal block is a real run}}
  \end{center}
  \vfill
\end{frame}

\end{document}
